\documentclass{beamer}
\usepackage[utf8]{inputenc}

\usetheme{Madrid}
\usecolortheme{default}
\usepackage{amsmath,amssymb,amsfonts,amsthm}
\usepackage{txfonts}
\usepackage{tkz-euclide}
\usepackage{listings}
\usepackage{adjustbox}
\usepackage{array}
\usepackage{tabularx}
\usepackage{gvv}
\usepackage{lmodern}
\usepackage{circuitikz}
\usepackage{tikz}
\usepackage{graphicx}

\setbeamertemplate{page number in head/foot}[totalframenumber]

\usepackage{tcolorbox}
\tcbuselibrary{minted,breakable,xparse,skins}



\definecolor{bg}{gray}{0.95}
\DeclareTCBListing{mintedbox}{O{}m!O{}}{%
  breakable=true,
  listing engine=minted,
  listing only,
  minted language=#2,
  minted style=default,
  minted options={%
    linenos,
    gobble=0,
    breaklines=true,
    breakafter=,,
    fontsize=\small,
    numbersep=8pt,
    #1},
  boxsep=0pt,
  left skip=0pt,
  right skip=0pt,
  left=25pt,
  right=0pt,
  top=3pt,
  bottom=3pt,
  arc=5pt,
  leftrule=0pt,
  rightrule=0pt,
  bottomrule=2pt,
  toprule=2pt,
  colback=bg,
  colframe=orange!70,
  enhanced,
  overlay={%
    \begin{tcbclipinterior}
    \fill[orange!20!white] (frame.south west) rectangle ([xshift=20pt]frame.north west);
    \end{tcbclipinterior}},
  #3,
}
\lstset{
    language=C,
    basicstyle=\ttfamily\small,
    keywordstyle=\color{blue},
    stringstyle=\color{orange},
    commentstyle=\color{green!60!black},
    numbers=left,
    numberstyle=\tiny\color{gray},
    breaklines=true,
    showstringspaces=false,
}
%------------------------------------------------------------

\title
{12.112}
\author 
{AI25BTECH11034 - Sujal Chauhan }



\begin{document}

\frame{\titlepage}
\begin{frame}{Question}
\textbf{Q.} The matrix
\[
A = \begin{pmatrix}
\frac{3}{2} & 0 & \frac{1}{2} \\
0 & -1 & 0 \\
\frac{1}{2} & 0 & \frac{3}{2}
\end{pmatrix}
\]
has three distinct eigenvalues and one of its eigenvectors is
\(
\begin{pmatrix}
1 \\ 0 \\ 1
\end{pmatrix}
\).
Which one of the following can be another eigenvector of $A$?

\[
\begin{array}{llll}
\text{(a)}~\begin{pmatrix} 1 \\ 0 \\ -1 \end{pmatrix} \quad &
\text{(b)}~\begin{pmatrix} -1 \\ 0 \\ 0 \end{pmatrix} \quad &
\text{(c)}~\begin{pmatrix} 0 \\ 1 \\ -1 \end{pmatrix} \quad &
\text{(d)}~\begin{pmatrix} 0\\ 0 \\ 1 \end{pmatrix}
\end{array}
\]


\end{frame}

\begin{frame}{Theory}
    \section*{Proof: Eigenvectors of a symmetric matrix corresponding to distinct eigenvalues are perpendicular}

Let $A$ be a symmetric matrix ($A^T = A$). Suppose $\mathbf{a}$ and $\mathbf{b}$ are eigenvectors of $A$ corresponding to distinct eigenvalues $\lambda_1$ and $\lambda_2$ respectively. That is:
\[
A\mathbf{a} = \lambda_1 \mathbf{a}, \qquad A\mathbf{b} = \lambda_2 \mathbf{b}
\]

Consider the product $\mathbf{a}^T A \mathbf{b}$ two ways:
\begin{align*}
\mathbf{a}^T A \mathbf{b} &= \mathbf{a}^T (\lambda_2 \mathbf{b}) = \lambda_2 \mathbf{a}^T \mathbf{b} \\
\end{align*}

But since $A$ is symmetric, 
\[
\mathbf{a}^T A \mathbf{b} = (A^T \mathbf{a})^T \mathbf{b} = (A \mathbf{a})^T \mathbf{b} = (\lambda_1 \mathbf{a})^T \mathbf{b} = \lambda_1 (\mathbf{a}^T \mathbf{b})
\]
\end{frame}
\begin{frame}{Theory}
Equating the two results:
\[
\lambda_1 (\mathbf{a}^T \mathbf{b}) = \lambda_2 (\mathbf{a}^T \mathbf{b})
\]

Since $\lambda_1 \neq \lambda_2$, it follows that
\[
(\lambda_1 - \lambda_2)(\mathbf{a}^T \mathbf{b}) = 0
\]
\[
\implies \mathbf{a}^T \mathbf{b} = 0
\]

Thus, the eigenvectors $\mathbf{a}$ and $\mathbf{b}$ are perpendicular.

\end{frame}
\begin{frame}{Solution}
   Let
\[
A=\begin{pmatrix}
\frac{3}{2} & 0 & \frac{1}{2} \\[4pt]
0 & -1 & 0 \\[4pt]
\frac{1}{2} & 0 & \frac{3}{2}
\end{pmatrix},
\qquad
a=\begin{pmatrix}1\\[2pt]0\\[2pt]1\end{pmatrix}.
\]
Since \(A\) is symmetric, eigenvectors corresponding to distinct eigenvalues are orthogonal.
Compute the inner product \(a^T b\) for each choice \(b\):
\end{frame}
\begin{frame}{Solution}
\[
\begin{aligned}
\text{(a)}\quad b&=\begin{pmatrix}1\\[2pt]0\\[2pt]-1\end{pmatrix}:
&\quad a^T b &= 1\cdot1 + 0\cdot0 + 1\cdot(-1)=0,\\[6pt]
\text{(b)}\quad b&=\begin{pmatrix}-1\\[2pt]0\\[2pt]0\end{pmatrix}:
&\quad a^T b &= -1\neq0,\\[6pt]
\text{(c)}\quad b&=\begin{pmatrix}0\\[2pt]1\\[2pt]-1\end{pmatrix}:
&\quad a^T b &= -1\neq0,\\[6pt]
\text{(d)}\quad b&=\begin{pmatrix}0\\[2pt]0\\[2pt]1\end{pmatrix}:
&\quad a^T b &= 1\neq0.
\end{aligned}
\]
\end{frame}
\begin{frame}{Solution}
Only (a) is orthogonal to \(a\); we now verify that it is indeed an eigenvector by direct multiplication:
\[
A\begin{pmatrix}1\\[2pt]0\\[2pt]-1\end{pmatrix}
=
\begin{pmatrix}
\frac32\cdot1 + \frac12\cdot(-1)\\[6pt]
0\\[6pt]
\frac12\cdot1 + \frac32\cdot(-1)
\end{pmatrix}
=
\begin{pmatrix}1\\[2pt]0\\[2pt]-1\end{pmatrix}.
\]
Hence \(\begin{pmatrix}1\\0\\-1\end{pmatrix}\) is an eigenvector (with eigenvalue \(1\)), so the correct choice is \(\boxed{\text{(a)}}\).




\end{frame}



\end{document}
